\section{Conclusion}
\label{sec:conclusion}

We began with a straightforward goal: measure how dataset documentation quality and benchmark concentration predict ML reproducibility failure, using HuggingFace Hub and OpenML as data sources. What we found instead was that the metadata infrastructure we assumed existed does not --- and that this absence follows a pattern that has implications for any researcher attempting programmatic reproducibility auditing of pre-2018 ML literature.

In this work, we designed and executed a systematic infrastructure feasibility study (H-E1) to determine whether HF Hub and OpenML can serve as data sources for a dataset-metadata-to-reproducibility-outcome regression study on Raff's 255-paper corpus. Our main contributions are:

\begin{enumerate}
  \item \textbf{The first empirical quantification of HF Hub and OpenML coverage for the Raff 2019 corpus:} HF Hub achieves 30\% dataset card coverage (15/50 datasets found, all NLP benchmarks); OpenML achieves 22\% pre-publication temporal coverage (11/50, all classical tabular). Both fall below the minimum thresholds required for valid regression analysis.

  \item \textbf{Characterization of the domain-temporal selection bias:} Coverage gaps precisely follow each platform's founding community. DL vision, speech, and graph datasets --- which dominate Raff's corpus --- are absent from both platforms.

  \item \textbf{Identification of Papers With Code as the appropriate primary data source:} Our findings provide a principled motivation for switching to Papers With Code API for a redesigned study.

  \item \textbf{A validated, reusable data acquisition pipeline:} RaffParser, HFCoverageChecker, OpenMLTemporalChecker, MetricsAggregator, and an 18-test suite are available for immediate reuse.
\end{enumerate}

\textbf{Future directions:} The highest-priority next experiment is a redesigned H-E1 using Papers With Code as primary API --- the most direct path to answering the original research question. Additional paths include systematic HF Hub name variant exploration (a lower-cost check of whether the 30\% figure is a lower bound) and extension to post-2019 corpora where HF coverage is substantially higher.

The infrastructure gap we discovered is both a limitation of this paper and a contribution to the field. Any researcher planning to use HF Hub and OpenML to study pre-2018 ML reproducibility should quantify coverage for their specific corpus before proceeding --- the structural domain biases in these platforms are invisible unless explicitly measured. We hope this work saves others from the same discovery at a later and more costly stage of their pipeline.
